\section{本讲小结}
动量定理把同一对象在同一时间区间内的合外冲量与动量变化联系起来。运算时先画速度和冲量方向，再写分量式；平均冲击力来自冲量除以作用时间。多阶段问题可利用全过程消去中间量，但须检查摩擦方向与事件顺序。

连续流问题要同时确定每单位质量的动量变化和每秒流过的质量；流量由边界与物质的相对速度确定。绳突然绷紧时，个体冲量方程与沿绳速度约束共同决定末速度，外部槽壁可能提供不可忽略的冲量。动能减少是求解后的物理核验，不能把理想绳误解为必然无耗散。
